% Front matter — tese curta, PT-PT.

\frontmatter
\pagestyle{plain}

\maketitlepage

%----------------------------------------------------------------------------------------
%	DECLARAÇÃO DE INTEGRIDADE
%----------------------------------------------------------------------------------------

\begin{soi}

\vspace{1cm}

Declaro ter conduzido este trabalho académico com integridade.

Não plagiei ou apliquei qualquer forma de uso indevido de informações ou falsificação de
resultados ao longo do processo que levou à sua elaboração.

Portanto, o trabalho apresentado neste documento é original e de minha autoria, não tendo sido
utilizado anteriormente para nenhum outro fim. As exceções estão explicitamente reconhecidas na
secção onde são abordadas as considerações éticas, a Secção~\ref{sec:met_etica}. A
Secção~\ref{sec:met_ia}, que se lhe segue, declara como as ferramentas de Inteligência Artificial
foram utilizadas e para que finalidade.

Declaro ainda que tenho pleno conhecimento do Código de Conduta Ética do P.PORTO.

\vspace{1cm}

% A redação acima é a do modelo oficial MEIA (v2, 2026), palavra por palavra, incluindo o período
% das exceções — que se aplica e por isso fica. É o modelo que manda a declaração de uso de IA
% para a secção das considerações éticas (Secção~\ref{sec:met_ia}), e não para esta página; a
% versão anterior deste ficheiro trazia-a aqui, o que contrariava o modelo.
% A data acompanha a da capa (definida em main.tex). Estava em \today e mudava sozinha a cada
% compilação, o que numa DECLARAÇÃO ASSINADA é pior do que na capa.
ISEP, Porto, setembro de 2026

\end{soi}

%----------------------------------------------------------------------------------------
%	DEDICATÓRIA
%----------------------------------------------------------------------------------------
\begin{dedicatory}
À minha família.
\end{dedicatory}

%----------------------------------------------------------------------------------------
%	RESUMO
%----------------------------------------------------------------------------------------

\begin{abstract}
Quando o preço de uma ação se move de forma abrupta, o investidor particular necessita de
determinar se o movimento é invulgar, se a sua origem é a empresa ou o mercado, e se existem
precedentes. As aplicações gratuitas apresentam a variação percentual, e as poucas que vão
além disso não expõem a evidência que as sustenta; os terminais profissionais respondem, e
não estão acessíveis a este investidor.

Esta dissertação apresenta o InvestiGator, um sistema de alertas financeiros assente exclusivamente
em fontes de dados gratuitas, que responde às três questões e expõe a evidência de cada
resposta. O sistema integra quatro técnicas: deteção de anomalias por \emph{z}-score deslizante,
decomposição do retorno em mercado, setor e empresa com encolhimento do beta, recuperação de
precedentes por semelhança semântica com medição do desfecho observado, e um modelo supervisionado
de triagem treinado sobre um conjunto de dados próprio.

Cada componente foi avaliado contra linhas de base transparentes. A deteção normalizada pela
volatilidade revela-se consistente entre empresas; a recuperação semântica supera a taxa-base nos
cinco setores avaliados; o modelo de triagem não supera uma linha de base apenas com
volatilidade, resultado negativo reportado e analisado.

Por desenho, o sistema não produz previsões de preço.
\end{abstract}

\begin{abstractotherlanguage}
When a stock price moves abruptly, retail investors need to determine whether the movement is
unusual, whether its origin lies with the company or the market, and whether comparable precedents
exist. Free applications display the percentage change, and the few that go further do not
expose the evidence behind them;
professional terminals do answer them, and are not accessible to this investor profile.

This dissertation presents InvestiGator, a financial alerting system built exclusively on free data
sources, which answers the three questions and exposes the evidence supporting each answer. The
system integrates four techniques: anomaly detection through a rolling \emph{z}-score, decomposition
of returns into market, sector and company components with beta shrinkage, precedent retrieval by
semantic similarity with measurement of the observed outcome, and a supervised triage model trained
on a purpose-built dataset.

Each component was evaluated against transparent baselines. Volatility-normalised detection proves
consistent across companies; semantic retrieval outperforms the base rate in every sector assessed;
the triage model does not outperform a volatility-only baseline, a negative result that is reported
and analysed.

By design, the system does not produce price forecasts.
\end{abstractotherlanguage}
%----------------------------------------------------------------------------------------
%	AGRADECIMENTOS
%----------------------------------------------------------------------------------------
% RASCUNHO NA VOZ DO AUTOR, PARA ELE REESCREVER. Os agradecimentos sao a unica
% seccao do documento que ninguem pode escrever por outra pessoa: o que aqui esta e uma
% base factual, com as pessoas certas e sem sentimentalismo, e nao a formulacao final.
%
% Duas regras que valeram na redaccao e que convem nao perder ao reescrever:
%   1. AGRADECER A QUEM PARTICIPOU, E SO A ESSES. Quem vota nos alertas do canal
%      participou mesmo, e o retorno esta reportado na Seccao 5.6.5. O que NAO se pode
%      escrever e' um agradecimento a participantes de um estudo controlado, porque esse
%      nao foi realizado e o Capitulo 6 declara-o.
%   2. A ESPECIFICIDADE SUBSTITUI A EMOCAO. "Pelas conversas que definiram a direccao"
%      diz mais, e soa melhor, do que "por todo o apoio".
\begin{acknowledgements}
Ao Prof.\ Luís Gomes, orientador, e a Rafael Silva, coorientador, pelas conversas que
definiram a direção deste trabalho, pelos exemplos que me deram a ver o que já existe na
área, e pela orientação na redação do documento. Parte das decisões aqui descritas foi
tomada depois de uma dessas conversas.

À minha família, pelo apoio ao longo destes meses.

Aos colegas da Sistrade, e à Sistrade, pela flexibilidade que tornou possível conciliar
este mestrado com o trabalho, e pelas conversas sobre a aplicação ao longo da sua
construção.

A quem acompanha o canal e classifica os alertas que recebe. Esse retorno é a única
informação deste trabalho que não veio de uma medição automática.
\end{acknowledgements}

%----------------------------------------------------------------------------------------
%	ÍNDICES
%----------------------------------------------------------------------------------------
\tableofcontents
% As listas passam a constar do Índice. A Lista de Acrónimos já constava e as outras
% três não, o que é a pior das configurações: o leitor aprende que as listas estão no
% Índice e depois não encontra lá as que procura. A convenção seguida é a da dissertação
% aprovada do Bruno Ribeiro, cujo Índice lista todas com página.
% O cleardoublepage antes do addcontentsline garante que a entrada aponta para a página
% onde a lista COMEÇA, e não para a anterior.
\cleardoublepage
\addcontentsline{toc}{chapter}{Lista de Figuras}
\listoffigures
\cleardoublepage
\addcontentsline{toc}{chapter}{Lista de Tabelas}
\listoftables
% A Lista de Excertos NAO e impressa e o comando de renomear o prefixo foi retirado: a
% reescrita para esta arvore deixou o documento SEM UM UNICO EXCERTO de codigo, e o
% comentario que aqui estava defendia manter uma lista que ja nao tem conteudo. A
% garantia anti-lookahead, que era o excerto que valia a pena mostrar, passou a ser a
% figura da janela deslizante nos metodos, que serve melhor o leitor a quem esta
% dissertacao se dirige. Verificado a 2026-09-04: zero ambientes lstlisting no corpo.

%----------------------------------------------------------------------------------------
%	LISTA DE SÍMBOLOS
%----------------------------------------------------------------------------------------
% Exigida pela lista de verificação do modelo oficial MEIA (v2, 2026), que a marca como
% obrigatória. Só entram os símbolos que o documento usa e define; a terceira coluna remete para
% a equação ou secção onde cada um é introduzido, para que a lista sirva de índice e não apenas
% de glossário. A Lista de Código e a Lista de Algoritmos não constam porque o documento não tem
% nenhum dos dois, e o próprio modelo manda remover esses comandos nesse caso.
\cleardoublepage
\addcontentsline{toc}{chapter}{Lista de Símbolos}
\begin{symbols}{l p{0.60\textwidth} l}
$P_t$ & Preço de fecho da sessão $t$ & Eq.~\ref{eq:zscore} \\
$r_t$ & Retorno logarítmico do dia $t$, $\ln(P_t / P_{t-1})$ & Eq.~\ref{eq:zscore} \\
$\mu$, $\sigma$ & Média e desvio-padrão dos retornos da janela anterior ao dia & Eq.~\ref{eq:zscore} \\
$z_t$ & Raridade do retorno do dia $t$ face à janela anterior & Eq.~\ref{eq:zscore} \\
\addlinespace
$r_m$, $\tilde{r}_s$ & Retorno diário do índice de mercado e do índice de setor & Eq.~\ref{eq:decomposicao} \\
$\beta_m$, $\beta_s$ & Sensibilidade histórica da ação ao mercado e ao setor & Eq.~\ref{eq:decomposicao} \\
$\varepsilon$ & Parte do retorno que não é explicada nem pelo mercado nem pelo setor & Eq.~\ref{eq:decomposicao} \\
$\hat{\beta}$ & Estimativa do beta obtida por regressão sobre a janela & Eq.~\ref{eq:vasicek} \\
$\beta_{\text{ref}}$ & Beta de referência para onde a estimativa é encolhida & Eq.~\ref{eq:vasicek} \\
$\mathrm{SE}(\hat{\beta})$ & Erro-padrão da estimativa do beta & Eq.~\ref{eq:vasicek} \\
$\sigma^2_{\text{ref}}$ & Variância de referência que fixa a força do encolhimento & Eq.~\ref{eq:vasicek} \\
$w$ & Peso do encolhimento de Vasicek, entre zero e um & Eq.~\ref{eq:vasicek} \\
\addlinespace
$\cos$ & Semelhança do cosseno entre dois vetores de frase & Secção~\ref{sec:met_recuperacao} \\
$h$ & Meia-vida do decaimento por recência aplicado aos precedentes & Secção~\ref{sec:met_recuperacao} \\
$k$ & Número de precedentes recuperados e apresentados & Secção~\ref{sec:met_recuperacao} \\
\addlinespace
$p_{\text{bruta}}$ & Pontuação de triagem antes da calibração & Eq.~\ref{eq:platt} \\
$p_{\text{calibrada}}$ & Pontuação de triagem depois da calibração de Platt & Eq.~\ref{eq:platt} \\
$a$, $b$ & Parâmetros aprendidos da calibração de Platt & Eq.~\ref{eq:platt} \\
\addlinespace
$P$, $C$ & Precisão e cobertura & Secção~\ref{sec:av_desenho} \\
$F_1$ & Média harmónica da precisão e da cobertura, \mbox{$2PC/(P+C)$} & Secção~\ref{sec:av_desenho} \\
\end{symbols}

% nogroupskip: sem isto o espaco vertical VARIA conforme a letra inicial -- entradas
% que partilham a inicial ficam coladas e as restantes separadas, o que se le como
% espacamento arbitrario.
\printnoidxglossary[type=\acronymtype, title=Lista de Acrónimos, nonumberlist,
                   nogroupskip]

\mainmatter
\pagestyle{thesis}
